% Câu hỏi mẫu: Vectơ trong không gian (Toán 12).
% Mỗi câu nằm trong \begin{ex}...\end{ex}; dòng "% @..." ngay phía trên khai báo thẻ, mức độ, mã câu.

% @tags: Lớp 12 / Vectơ trong không gian / Tích vô hướng ; @level: NB ; @id: 12V-001
\begin{ex}
  Trong không gian $Oxyz$, cho $\vec u=(1;2;-1)$ và $\vec v=(2;0;1)$. Tích vô hướng $\vec u\cdot\vec v$ bằng
  \choice
    {$0$}
    {\True $1$}
    {$2$}
    {$3$}
  \loigiai{$\vec u\cdot\vec v = 1\cdot 2 + 2\cdot 0 + (-1)\cdot 1 = 1$.}
\end{ex}

% @tags: Lớp 12 / Vectơ trong không gian / Tọa độ điểm ; @level: NB ; @id: 12V-002
\begin{ex}
  Trong không gian $Oxyz$, cho $A(1;2;3)$ và $B(3;0;-1)$. Tọa độ trung điểm $I$ của đoạn $AB$ là
  \choice
    {\True $(2;1;1)$}
    {$(4;2;2)$}
    {$(1;-1;-2)$}
    {$(2;-2;-4)$}
  \loigiai{$I\left(\dfrac{1+3}{2};\dfrac{2+0}{2};\dfrac{3+(-1)}{2}\right)=(2;1;1)$.}
\end{ex}

% Câu dưới đây chỉ có mã ID kiểu ex_test: hệ thống tự suy ra Lớp 12 và mức Nhận biết.
\begin{ex}%[2H3N1-1]
  Trong không gian $Oxyz$, độ dài của vectơ $\vec u=(2;-1;2)$ bằng
  \choice
    {\True $3$}
    {$9$}
    {$\sqrt5$}
    {$5$}
  \loigiai{$|\vec u|=\sqrt{2^2+(-1)^2+2^2}=\sqrt9=3$.}
\end{ex}

% @tags: Lớp 12 / Vectơ trong không gian / Tích vô hướng ; @level: TH ; @id: 12V-004
\begin{ex}
  Trong không gian $Oxyz$, cho $\vec a=(1;m;2)$ và $\vec b=(2;-1;1)$. Giá trị của $m$ để $\vec a\perp\vec b$ là
  \choice
    {\True $m=4$}
    {$m=-4$}
    {$m=2$}
    {$m=0$}
  \loigiai{$\vec a\perp\vec b \Leftrightarrow 2-m+2=0 \Leftrightarrow m=4$.}
\end{ex}

% @tags: Lớp 12 / Vectơ trong không gian / Tọa độ điểm ; @level: TH ; @id: 12V-005
\begin{ex}
  Trong không gian $Oxyz$, cho $A(1;0;0)$, $B(0;2;0)$, $C(0;0;3)$. Tọa độ trọng tâm $G$ của tam giác $ABC$ là
  \choice
    {$(1;2;3)$}
    {\True $\left(\dfrac13;\dfrac23;1\right)$}
    {$\left(\dfrac12;1;\dfrac32\right)$}
    {$(3;2;1)$}
  \loigiai{$G\left(\dfrac{1+0+0}{3};\dfrac{0+2+0}{3};\dfrac{0+0+3}{3}\right)=\left(\dfrac13;\dfrac23;1\right)$.}
\end{ex}

% @tags: Lớp 12 / Vectơ trong không gian / Tích vô hướng ; @level: TH ; @id: 12V-006
\begin{ex}
  Trong không gian $Oxyz$, góc giữa hai vectơ $\vec u=(1;0;1)$ và $\vec v=(0;1;1)$ bằng
  \choice
    {$30^\circ$}
    {$45^\circ$}
    {\True $60^\circ$}
    {$90^\circ$}
  \loigiai{$\cos(\vec u,\vec v)=\dfrac{0+0+1}{\sqrt2\cdot\sqrt2}=\dfrac12$ nên góc bằng $60^\circ$.}
\end{ex}

% Lựa chọn "Cả ba vectơ trên" được ghim (\Pin) nên luôn ở vị trí D khi đảo.
% @tags: Lớp 12 / Vectơ trong không gian / Tích vô hướng ; @level: TH ; @id: 12V-007
\begin{ex}
  Trong không gian $Oxyz$, vectơ nào sau đây vuông góc với $\vec u=(1;1;0)$?
  \choice
    {$\vec a=(1;-1;0)$}
    {$\vec b=(0;0;1)$}
    {$\vec c=(1;-1;5)$}
    {\True\Pin Cả ba vectơ trên}
  \loigiai{$\vec u\cdot\vec a=1-1+0=0$, $\vec u\cdot\vec b=0$, $\vec u\cdot\vec c=1-1+0=0$. Cả ba vectơ đều vuông góc với $\vec u$.}
\end{ex}

% @tags: Lớp 12 / Vectơ trong không gian / Tích có hướng ; @level: VD ; @id: 12V-008
\begin{ex}
  Trong không gian $Oxyz$, cho $A(1;1;1)$, $B(2;3;4)$, $C(4;1;3)$. Diện tích tam giác $ABC$ bằng
  \choice
    {\True $\dfrac{\sqrt{101}}{2}$}
    {$\sqrt{101}$}
    {$\dfrac{101}{2}$}
    {$\dfrac{\sqrt{101}}{3}$}
  \loigiai{$\overrightarrow{AB}=(1;2;3)$, $\overrightarrow{AC}=(3;0;2)$, $\left[\overrightarrow{AB},\overrightarrow{AC}\right]=(4;7;-6)$.
  Vậy $S_{ABC}=\dfrac12\sqrt{16+49+36}=\dfrac{\sqrt{101}}{2}$.}
\end{ex}

% @tags: Lớp 12 / Vectơ trong không gian / Tọa độ điểm ; @level: VD ; @id: 12V-009
\begin{ex}
  Trong không gian $Oxyz$, cho $A(1;0;2)$, $B(2;1;0)$, $C(0;3;1)$. Tọa độ điểm $D$ để $ABCD$ là hình bình hành là
  \choice
    {\True $(-1;2;3)$}
    {$(3;-2;1)$}
    {$(1;4;-1)$}
    {$(-1;-2;-3)$}
  \loigiai{$ABCD$ là hình bình hành $\Leftrightarrow \overrightarrow{AD}=\overrightarrow{BC}$, suy ra $D=A+C-B=(-1;2;3)$.}
\end{ex}

% @tags: Lớp 12 / Vectơ trong không gian / Tọa độ điểm ; @level: VD ; @id: 12V-010
\begin{ex}
  Trong không gian $Oxyz$, cho $A(1;2;3)$, $B(3;-2;1)$. Điểm $M$ thuộc trục $Ox$ sao cho $MA^2+MB^2$ nhỏ nhất có tọa độ là
  \choice
    {$(1;0;0)$}
    {\True $(2;0;0)$}
    {$(3;0;0)$}
    {$(4;0;0)$}
  \loigiai{Với $M(x;0;0)$: $MA^2+MB^2=(x-1)^2+13+(x-3)^2+5=2x^2-8x+28=2(x-2)^2+20$, nhỏ nhất khi $x=2$.}
\end{ex}

% @tags: Lớp 12 / Vectơ trong không gian / Tọa độ điểm ; @level: TH ; @id: 12V-011
\begin{ex}
  Trong không gian $Oxyz$, cho $A(1;0;0)$, $B(0;1;0)$, $C(0;0;1)$. Xét tính đúng sai của các mệnh đề sau.
  \choiceTF
    {\True $\overrightarrow{AB}=(-1;1;0)$}
    {\True Tam giác $ABC$ là tam giác đều}
    {$\overrightarrow{AB}\cdot\overrightarrow{AC}=0$}
    {\True Diện tích tam giác $ABC$ bằng $\dfrac{\sqrt3}{2}$}
  \loigiai{$AB=BC=CA=\sqrt2$ nên tam giác đều, diện tích $\dfrac{\sqrt3}{4}\cdot 2=\dfrac{\sqrt3}{2}$.
  $\overrightarrow{AB}\cdot\overrightarrow{AC}=(-1)(-1)+0+0=1\neq0$.}
\end{ex}

% @tags: Lớp 12 / Vectơ trong không gian / Tích vô hướng ; @level: TH ; @id: 12V-012
\begin{ex}
  Trong không gian $Oxyz$, cho $\vec u=(1;-2;2)$ và $\vec v=(2;1;0)$. Xét tính đúng sai của các mệnh đề sau.
  \choiceTF
    {\True $|\vec u|=3$}
    {\True $\vec u\perp\vec v$}
    {$\vec u-\vec v=(-1;-3;3)$}
    {\True Vectơ $2\vec u$ có độ dài bằng $6$}
  \loigiai{$|\vec u|=\sqrt{1+4+4}=3$; $\vec u\cdot\vec v=2-2+0=0$; $\vec u-\vec v=(-1;-3;2)$; $|2\vec u|=2|\vec u|=6$.}
\end{ex}

% @tags: Lớp 12 / Vectơ trong không gian / Tọa độ điểm ; @level: NB ; @id: 12V-013
\begin{ex}
  Trong không gian $Oxyz$, tính độ dài của vectơ $\vec u=(3;4;0)$.
  \shortans{5}
  \loigiai{$|\vec u|=\sqrt{9+16+0}=5$.}
\end{ex}

% @tags: Lớp 12 / Vectơ trong không gian / Tích vô hướng ; @level: TH ; @id: 12V-014
\begin{ex}
  Trong không gian $Oxyz$, cho $A(1;1;1)$, $B(2;3;4)$, $C(4;1;3)$. Tính $\overrightarrow{AB}\cdot\overrightarrow{AC}$.
  \shortans{9}
  \loigiai{$\overrightarrow{AB}=(1;2;3)$, $\overrightarrow{AC}=(3;0;2)$, nên $\overrightarrow{AB}\cdot\overrightarrow{AC}=3+0+6=9$.}
\end{ex}

% @tags: Lớp 12 / Vectơ trong không gian / Tọa độ điểm ; @level: VD ; @id: 12V-015
\begin{ex}
  Trong không gian $Oxyz$, cho $A(2;0;0)$, $B(0;2;0)$, $C(0;0;2)$.
  \begin{enumerate}[label=\alph*)]
    \item Chứng minh tam giác $ABC$ đều.
    \item Tìm tọa độ điểm $D$ sao cho $ABCD$ là hình bình hành.
  \end{enumerate}
  \loigiai{a) $AB=BC=CA=2\sqrt2$ nên tam giác $ABC$ đều.

  b) $\overrightarrow{AD}=\overrightarrow{BC}=(0;-2;2)$ nên $D=(2;-2;2)$.}
\end{ex}
